\begingroup\setbox0=\vbox{\section*{(*) Bài toán hai hiệu số}}\endgroup
\chapterhead{9}{(*) Bài toán hai hiệu số}

\shead{So sánh hai cách chia}
\begin{knowbox}
Nhiều bài toán mô tả \emph{cùng một lượng} được chia (hoặc xếp) theo \textbf{hai cách khác nhau}. Ta so sánh:
\begin{itemize}[leftmargin=1.4em,itemsep=3pt]
\item \textbf{Hiệu số phần mỗi người (mỗi đơn vị):} mỗi phần (mỗi em, mỗi ghế, mỗi giờ,\ldots) ở cách này nhiều hơn ở cách kia bao nhiêu.
\item \textbf{Hiệu số tổng:} cách này cần nhiều hơn cách kia bao nhiêu. Nếu một cách \emph{thừa} và cách kia \emph{thiếu} thì hiệu số tổng là \emph{số thừa $+$ số thiếu}; nếu cả hai cách cùng thừa (hoặc cùng thiếu) thì lấy số lớn trừ số bé.
\end{itemize}
\textbf{Số phần} $=$ (hiệu số tổng) $:$ (hiệu số mỗi phần). Từ đó tìm đại lượng cần tìm.\par
\smallskip
\emph{Phương pháp giả thiết tạm} cũng dựa trên ý tưởng này: tạm giả sử tất cả đều cùng một loại, rồi so sánh với thực tế.
\end{knowbox}
\gap
\begin{exbox}{\ExWord\ 1}
Cô giáo chia kẹo cho một nhóm học sinh. Nếu mỗi em $5$ cái thì thừa $8$ cái; nếu mỗi em $7$ cái thì thiếu $4$ cái. Hỏi có bao nhiêu em và bao nhiêu cái kẹo?
\guide
Lần thứ hai mỗi em được nhiều hơn $7-5=2$ cái.\par
Cách thứ hai cần nhiều hơn $8+4=12$ cái kẹo ($8$ cái thừa, thêm $4$ cái nữa).\par
Số em: $12:2=6$.\par Số kẹo: $5\times6+8=38$. \ (\emph{Thử lại:} $7\times6-4=38$.)
\dapso{$6$ em, $38$ cái kẹo.}
\end{exbox}
\gap
\begin{exbox}{\ExWord\ 2}
Lớp 5B kê ghế để họp. Nếu mỗi ghế $3$ học sinh thì $6$ học sinh không có chỗ ngồi. Nếu mỗi ghế $4$ học sinh thì thừa $2$ ghế. Hỏi có bao nhiêu ghế và bao nhiêu học sinh?
\guide
Thừa hai ghế ở cách thứ hai nghĩa là cần thêm $4\times2=8$ học sinh nữa mới ngồi kín tất cả các ghế.\par
Hiệu số mỗi ghế: $4-3=1$ học sinh.\par Hiệu số tổng: $6+8=14$ học sinh.\par
Số ghế: $14:1=14$.\par Số học sinh: $3\times14+6=48$.
\dapso{$14$ ghế, $48$ học sinh.}
\end{exbox}
\gap
\begin{exbox}{\ExWord\ 3}
Vừa gà vừa chó có tất cả $36$ con và $100$ chân. Hỏi có bao nhiêu con gà, bao nhiêu con chó?
\guide
\emph{Giả sử} cả $36$ con đều là gà thì có $36\times2=72$ chân, ít hơn thực tế $100-72=28$ chân.\par
Thay một con gà bằng một con chó thì số chân tăng thêm $4-2=2$ chân.\par Số chó: $28:2=14$; số gà: $36-14=22$.
\dapso{$22$ con gà, $14$ con chó.}
\end{exbox}
\gap
\begin{exbox}{\ExWord\ 4}
Một người đi từ nhà lên thị xã. Nếu mỗi giờ đi $12$ km thì đến muộn hơn dự định $1$ giờ; nếu mỗi giờ đi $15$ km thì đến sớm hơn $1$ giờ. Tìm quãng đường và thời gian dự định.
\guide
Trong đúng thời gian dự định, nếu đi $12$ km mỗi giờ thì còn thiếu $12$ km (giờ đến muộn); nếu đi $15$ km mỗi giờ thì vượt quá $15$ km.\par
Hiệu số mỗi giờ: $15-12=3$ km; hiệu số tổng: $12+15=27$ km.\par Thời gian dự định: $27:3=9$ giờ.\par
Quãng đường: $12\times(9+1)=120$ km.
\dapso{$120$ km; $9$ giờ.}
\end{exbox}
\dividerbar{Luyện tập}
\begin{probbox}{1}{\relax}
Chia vở: mỗi em $4$ quyển thì thừa $5$ quyển; mỗi em $5$ quyển thì thiếu $3$ quyển. Tìm số học sinh và số vở.
\end{probbox}
\probsol{
Hiệu số mỗi em: $1$; hiệu số tổng: $5+3=8$.\par Số học sinh: $8$; số vở: $4\times8+5=37$.
\dapso{$8$ học sinh, $37$ quyển vở.}}
\gap
\begin{probbox}{2}{\relax}
Chia bánh: mỗi em $3$ cái thì thừa $16$ cái; mỗi em $5$ cái thì vẫn thừa $4$ cái. Hỏi có bao nhiêu em và bao nhiêu cái bánh?
\end{probbox}
\probsol{
Cả hai cách đều thừa: hiệu số tổng $16-4=12$; mỗi em $2$.\par Số em: $6$; số bánh: $3\times6+16=34$.
\dapso{$6$ em, $34$ cái bánh.}}
\gap
\begin{probbox}{3}{\relax}
Chia cam: mỗi em $6$ quả thì thiếu $2$ quả; mỗi em $8$ quả thì thiếu $14$ quả. Tìm số em và số cam.
\end{probbox}
\probsol{
Cả hai cách đều thiếu: hiệu số tổng $14-2=12$; mỗi em $2$.\par Số em: $6$; số cam: $6\times6-2=34$.
\dapso{$6$ em, $34$ quả cam.}}
\gap
\begin{probbox}{4}{\relax}
Một trường thuê xe ô tô đi tham quan. Nếu mỗi xe chở $40$ học sinh thì $10$ em không có chỗ; nếu mỗi xe chở $45$ học sinh thì thừa $15$ chỗ. Hỏi có bao nhiêu xe và bao nhiêu học sinh?
\end{probbox}
\probsol{
Mỗi xe: $5$ chỗ; tổng cộng: $10+15=25$.\par Số xe: $25:5=5$.\par Số học sinh: $40\times5+10=210$.
\dapso{$5$ xe, $210$ học sinh.}}
\gap
\begin{probbox}{5}{\relax}
Mẹ mang một số tiền đi chợ. Nếu mua $5$ kg táo thì thiếu $20\,000$ đồng; nếu mua $4$ kg thì thừa $10\,000$ đồng. Tìm giá $1$ kg táo và số tiền mẹ mang đi.
\end{probbox}
\probsol{
Mua thêm một ki-lô-gam thì chênh lệch $20\,000+10\,000=30\,000$ đồng, nên $1$ kg giá $30\,000$ đồng.\par Số tiền: $30\,000\times4+10\,000=130\,000$ đồng.
\dapso{$30\,000$ đồng; $130\,000$ đồng.}}
\gap
\begin{probbox}{6}{\relax}
Nếu mỗi ghế ngồi $5$ người thì $4$ người phải đứng; nếu mỗi ghế ngồi $6$ người thì thừa $1$ ghế. Hỏi có bao nhiêu ghế và bao nhiêu người?
\end{probbox}
\probsol{
Thừa một ghế nghĩa là thiếu $6$ người.\par Hiệu số tổng: $4+6=10$; mỗi ghế $1$.\par Số ghế: $10$; số người: $5\times10+4=54$.
\dapso{$10$ ghế, $54$ người.}}
\gap
\begin{probbox}{7}{\relax}
Một đoàn khách đến nhà nghỉ. Nếu mỗi phòng $4$ người thì $6$ khách không có phòng; nếu mỗi phòng $6$ người thì thừa $2$ phòng. Hỏi có bao nhiêu phòng và bao nhiêu khách?
\end{probbox}
\probsol{
Thừa hai phòng nghĩa là thiếu $12$ khách.\par Hiệu số tổng $6+12=18$; mỗi phòng $2$.\par Số phòng: $9$; số khách: $4\times9+6=42$.
\dapso{$9$ phòng, $42$ khách.}}
\gap
\begin{probbox}{8}{\relax}
Tìm một số, biết rằng $7$ lần số đó vượt quá $13$ đơn vị so với $50$, còn $5$ lần số đó thì kém $5$ đơn vị so với $50$.
\end{probbox}
\probsol{
Hai tích hơn kém nhau $13+5=18$, hai thừa số hơn kém nhau $7-5=2$.\par Số cần tìm là $18:2=9$.
\dapso{$9$.}}
\gap
\begin{probbox}{9}{\relax}
Mua $3$ bút và $5$ quyển vở hết $49\,000$ đồng; mua $3$ bút và $8$ quyển vở cùng loại hết $64\,000$ đồng. Tìm giá một cái bút và một quyển vở.
\end{probbox}
\probsol{
Hai lần mua chênh nhau $3$ quyển vở và $64\,000-49\,000=15\,000$ đồng, nên một quyển vở giá $5\,000$ đồng.\par
$3$ cái bút giá $49\,000-25\,000=24\,000$ đồng, nên một cái bút giá $8\,000$ đồng.
\dapso{Bút: $8\,000$ đồng; vở: $5\,000$ đồng.}}
\gap
\begin{probbox}{10}{\relax}
Mua $4$ bút chì và $3$ cục tẩy hết $22\,000$ đồng; mua $2$ bút chì và $3$ cục tẩy hết $14\,000$ đồng. Tìm giá mỗi loại.
\end{probbox}
\probsol{
$2$ bút chì giá $22\,000-14\,000=8\,000$ đồng, nên một bút chì giá $4\,000$ đồng.\par $3$ cục tẩy giá $14\,000-8\,000=6\,000$ đồng, nên một cục giá $2\,000$ đồng.
\dapso{Bút chì: $4\,000$ đồng; tẩy: $2\,000$ đồng.}}
\gap
\begin{probbox}{11}{\relax}
$30$ học sinh lớp 5C mua $100$ quyển vở; mỗi bạn nam mua $4$ quyển, mỗi bạn nữ mua $3$ quyển. Hỏi có bao nhiêu bạn nam, bao nhiêu bạn nữ?
\end{probbox}
\probsol{
Nếu cả $30$ bạn đều là nữ thì mua $90$ quyển, thiếu $10$ quyển.\par Mỗi bạn nam mua nhiều hơn mỗi bạn nữ $1$ quyển, nên có $10$ bạn nam và $20$ bạn nữ.
\dapso{$10$ bạn nam, $20$ bạn nữ.}}
\gap
\begin{probbox}{12}{\relax}
Một người đi xe đạp từ A đến B. Nếu mỗi giờ đi $10$ km thì đến muộn $1$ giờ; nếu mỗi giờ đi $15$ km thì đến sớm $1$ giờ. Tìm quãng đường AB.
\end{probbox}
\probsol{
Trong thời gian dự định, đi $10$ km mỗi giờ thì còn thiếu $10$ km; đi $15$ km mỗi giờ thì vượt quá $15$ km.\par Thời gian dự định: $(10+15):(15-10)=5$ giờ.\par
Quãng đường: $10\times(5+1)=60$ km.
\dapso{$60$ km.}}
\gap
\begin{probbox}{13}{\relax}
Chia kẹo cho một nhóm: mỗi em $4$ cái thì thừa $6$ cái; mỗi em $6$ cái thì thiếu $4$ cái. Hỏi có bao nhiêu em và bao nhiêu cái kẹo?
\end{probbox}
\probsol{
Hai cách chia chênh nhau $6+4=10$ cái kẹo.\par Mỗi em chênh nhau $6-4=2$ cái.\par Số em: $10:2=5$.\par Số kẹo: $4\times5+6=26$.
\dapso{$5$ em; $26$ cái kẹo.}}
\gap
\begin{probbox}{14}{\relax}
Chia vở: mỗi em $3$ quyển thì thừa $12$ quyển; mỗi em $5$ quyển thì vẫn thừa $2$ quyển. Tìm số học sinh và số vở.
\end{probbox}
\probsol{
Số vở thừa chênh nhau $12-2=10$.\par Mỗi em chênh nhau $5-3=2$.\par Số học sinh: $10:2=5$.\par Số vở: $3\times5+12=27$.
\dapso{$5$ học sinh; $27$ quyển vở.}}
\gap
\begin{probbox}{15}{\relax}
Chia quýt: mỗi em $7$ quả thì thiếu $3$ quả; mỗi em $9$ quả thì thiếu $15$ quả. Tìm số em và số quýt.
\end{probbox}
\probsol{
Số quýt thiếu chênh nhau $15-3=12$.\par Mỗi em chênh nhau $9-7=2$.\par Số em: $12:2=6$.\par Số quýt: $7\times6-3=39$.
\dapso{$6$ em; $39$ quả quýt.}}
\gap
\begin{probbox}{16}{\relax}
Vừa gà vừa thỏ có $30$ con, tất cả có $84$ chân. Hỏi mỗi loại có bao nhiêu con?
\end{probbox}
\probsol{
Giả sử cả $30$ con đều là gà: $2\times30=60$ chân.\par Số chân còn thiếu: $84-60=24$.\par Mỗi con thỏ nhiều hơn $2$ chân, nên số thỏ: $24:2=12$.\par Số gà: $30-12=18$.
\dapso{$18$ con gà, $12$ con thỏ.}}
\gap
\begin{probbox}{17}{\relax}
Một bãi giữ xe có $25$ chiếc xe đạp ($2$ bánh) và xích lô ($3$ bánh), tất cả có $62$ bánh xe. Hỏi mỗi loại có bao nhiêu chiếc?
\end{probbox}
\probsol{
Giả sử cả $25$ chiếc đều là xe đạp: $2\times25=50$ bánh.\par Còn thiếu: $62-50=12$ bánh.\par Mỗi chiếc xích lô nhiều hơn $1$ bánh, nên số xích lô: $12$.\par Số xe đạp: $25-12=13$.
\dapso{$13$ xe đạp, $12$ xích lô.}}
\gap
\begin{probbox}{18}{\relax}
Lan có $20$ tờ tiền loại $2\,000$ đồng và $5\,000$ đồng, tổng cộng $64\,000$ đồng. Hỏi mỗi loại có bao nhiêu tờ?
\end{probbox}
\probsol{
Giả sử tất cả đều là tờ $2\,000$ đồng: $2\,000\times20=40\,000$ đồng.\par Còn thiếu: $64\,000-40\,000=24\,000$ đồng.\par Mỗi tờ $5\,000$ đồng hơn $3\,000$ đồng.\par Số tờ $5\,000$ đồng: $24\,000:3\,000=8$.\par Số tờ $2\,000$ đồng: $20-8=12$.
\dapso{$12$ tờ $2\,000$ đồng, $8$ tờ $5\,000$ đồng.}}
\gap
\begin{probbox}{19}{\relax}
Mua $2$ bút và $3$ quyển vở hết $34\,000$ đồng; mua $2$ bút và $5$ quyển vở hết $46\,000$ đồng. Tìm giá mỗi loại.
\end{probbox}
\probsol{
Lần thứ hai mua nhiều hơn $2$ quyển vở và trả nhiều hơn $12\,000$ đồng.\par Một quyển vở: $12\,000:2=6\,000$ đồng.\par $2$ cái bút: $34\,000-6\,000\times3=16\,000$ đồng.\par Một cái bút: $8\,000$ đồng.
\dapso{Bút $8\,000$ đồng, vở $6\,000$ đồng.}}
\gap
\begin{probbox}{20}{\relax}
Mua $3$ kg cam và $2$ kg táo hết $150\,000$ đồng; mua $5$ kg cam và $2$ kg táo hết $200\,000$ đồng. Tìm giá một ki-lô-gam mỗi loại.
\end{probbox}
\probsol{
$2$ kg cam giá $200\,000-150\,000=50\,000$ đồng.\par $1$ kg cam: $25\,000$ đồng.\par $2$ kg táo: $150\,000-25\,000\times3=75\,000$ đồng.\par $1$ kg táo: $37\,500$ đồng.
\dapso{Cam $25\,000$ đồng/kg, táo $37\,500$ đồng/kg.}}
\gap
\begin{probbox}{21}{\relax}
Tìm một số, biết rằng $6$ lần số đó vượt quá $14$ đơn vị so với $40$, còn $4$ lần số đó thì kém $4$ đơn vị so với $40$.
\end{probbox}
\probsol{
Hai tích hơn kém nhau $14+4=18$.\par Hai thừa số hơn kém nhau $6-4=2$.\par Số cần tìm: $18:2=9$.\par Thử lại: $9\times6=54=40+14$; $9\times4=36=40-4$.
\dapso{$9$.}}
\gap
\begin{probbox}{22}{\relax}
Nếu mỗi ghế ngồi $4$ người thì $3$ người phải đứng; nếu mỗi ghế ngồi $5$ người thì thừa $2$ ghế. Hỏi có bao nhiêu ghế và bao nhiêu người?
\end{probbox}
\probsol{
$2$ ghế thừa có thể ngồi được $5\times2=10$ người.\par Hai cách xếp chênh nhau $3+10=13$ người.\par Mỗi ghế chênh nhau $5-4=1$.\par Số ghế: $13$.\par Số người: $4\times13+3=55$.
\dapso{$13$ ghế, $55$ người.}}
\gap
\begin{probbox}{23}{\relax}
Học sinh xếp hàng. Nếu mỗi hàng $6$ em thì thừa $4$ em; nếu mỗi hàng $7$ em (số hàng không đổi) thì thiếu $3$ em. Hỏi có bao nhiêu hàng và bao nhiêu học sinh?
\end{probbox}
\probsol{
Hai cách xếp chênh nhau $4+3=7$ học sinh.\par Mỗi hàng chênh nhau $1$.\par Số hàng: $7$.\par Số học sinh: $6\times7+4=46$.
\dapso{$7$ hàng, $46$ học sinh.}}
\gap
\begin{probbox}{24}{\relax}
Một người đi từ A đến B. Nếu mỗi giờ đi $15$ km thì đến muộn $2$ giờ; nếu mỗi giờ đi $20$ km thì đến sớm $1$ giờ. Tìm thời gian dự định và quãng đường AB.
\end{probbox}
\probsol{
Đi $15$ km/giờ trong thời gian dự định thì còn thiếu $15\times2=30$ km.\par Đi $20$ km/giờ trong thời gian dự định thì vượt quá $20\times1=20$ km.\par Hiệu số là $30+20=50$ km, mỗi giờ chênh $20-15=5$ km.\par Thời gian dự định: $50:5=10$ giờ.\par AB: $15\times(10+2)=180$ km.
\dapso{$10$ giờ; $180$ km.}}
\gap
\begin{probbox}{25}{\relax}
Một bài kiểm tra có $10$ câu: mỗi câu đúng được $10$ điểm, mỗi câu sai bị trừ $5$ điểm. Minh trả lời hết các câu và được $70$ điểm. Hỏi Minh trả lời đúng bao nhiêu câu?
\end{probbox}
\probsol{
Nếu đúng tất cả thì được $10\times10=100$ điểm.\par Minh bị thiếu $100-70=30$ điểm.\par Mỗi câu sai làm mất $10+5=15$ điểm.\par Số câu sai: $30:15=2$; số câu đúng: $10-2=8$.
\dapso{$8$ câu.}}
\gap
\begin{probbox}{26}{\relax}
$20$ con nhện ($8$ chân) và bọ cánh cứng ($6$ chân) có tất cả $136$ chân. Hỏi mỗi loại có bao nhiêu con?
\end{probbox}
\probsol{
Giả sử tất cả đều là bọ cánh cứng: $6\times20=120$ chân.\par Còn thiếu: $136-120=16$.\par Mỗi con nhện nhiều hơn $2$ chân, nên số nhện: $16:2=8$.\par Số bọ cánh cứng: $20-8=12$.
\dapso{$8$ con nhện, $12$ con bọ cánh cứng.}}
\gap
\begin{probbox}{27}{\relax}
$10$ chiếc thuyền, thuyền lớn chở $6$ người, thuyền nhỏ chở $4$ người, chở vừa đủ $48$ người. Hỏi mỗi loại có bao nhiêu chiếc?
\end{probbox}
\probsol{
Giả sử tất cả đều là thuyền nhỏ: $4\times10=40$ người.\par Còn thiếu: $48-40=8$.\par Mỗi thuyền lớn chở nhiều hơn $2$ người, nên số thuyền lớn: $8:2=4$.\par Số thuyền nhỏ: $10-4=6$.
\dapso{$4$ thuyền lớn, $6$ thuyền nhỏ.}}
\gap
\begin{probbox}{28}{\relax}
Bình mua $16$ con tem loại $3\,000$ đồng và $5\,000$ đồng, trả hết $62\,000$ đồng. Hỏi Bình mua mỗi loại bao nhiêu con tem?
\end{probbox}
\probsol{
Giả sử tất cả đều loại $3\,000$ đồng: $3\,000\times16=48\,000$ đồng.\par Còn thiếu: $62\,000-48\,000=14\,000$ đồng.\par Mỗi con tem $5\,000$ đồng đắt hơn $2\,000$ đồng.\par Số tem $5\,000$ đồng: $14\,000:2\,000=7$.\par Số tem $3\,000$ đồng: $16-7=9$.
\dapso{$9$ con tem loại $3\,000$ đồng, $7$ con tem loại $5\,000$ đồng.}}
\gap
\begin{probbox}{29}{\relax}
Mua $4$ quyển vở và $2$ cái bút hết $40\,000$ đồng; mua $2$ quyển vở và $3$ cái bút hết $36\,000$ đồng. Tìm giá mỗi loại.
\end{probbox}
\probsol{
Gấp đôi lần mua thứ hai: $4$ quyển vở và $6$ cái bút giá $72\,000$ đồng.\par So với lần mua thứ nhất, nhiều hơn $4$ cái bút và nhiều hơn $32\,000$ đồng.\par Một cái bút: $32\,000:4=8\,000$ đồng.\par Một quyển vở: $(40\,000-8\,000\times2):4=6\,000$ đồng.
\dapso{Vở $6\,000$ đồng, bút $8\,000$ đồng.}}
\gap
\begin{probbox}{30}{\relax}
Chia bóng cho các lớp: mỗi lớp $5$ quả thì thừa $3$ quả; mỗi lớp $7$ quả thì thiếu $5$ quả. Tìm số lớp và số bóng.
\end{probbox}
\probsol{
Hai cách chia chênh nhau $3+5=8$ quả bóng.\par Mỗi lớp chênh nhau $2$.\par Số lớp: $8:2=4$.\par Số bóng: $5\times4+3=23$.
\dapso{$4$ lớp, $23$ quả bóng.}}
\dividerbar{Thử thách -- Nâng cao}
\begin{probbox}{31}{\stars{1}}
Chia kẹo: mỗi em $3$ cái thì thừa $2$ cái; mỗi em $4$ cái thì có một em chỉ được $2$ cái. Hỏi có bao nhiêu em và bao nhiêu cái kẹo?
\end{probbox}
\probsol{
Cách thứ hai thiếu $4-2=2$ cái.\par Hiệu số tổng $2+2=4$; mỗi em $1$.\par Số em: $4$; số kẹo: $3\times4+2=14$.
\dapso{$4$ em, $14$ cái kẹo.}}
\gap
\begin{probbox}{32}{\stars{1}}
Học sinh xếp hàng. Nếu mỗi hàng $8$ em thì thừa $3$ em; nếu mỗi hàng $9$ em (số hàng không đổi) thì hàng cuối thiếu $5$ em. Hỏi có bao nhiêu học sinh?
\end{probbox}
\probsol{
Hiệu số tổng $3+5=8$; mỗi hàng $1$.\par Số hàng: $8$; số học sinh: $8\times8+3=67$.
\dapso{$67$ học sinh.}}
\gap
\begin{probbox}{33}{\stars{1}}
Một tổ dự định mỗi ngày làm $30$ sản phẩm. Thực tế mỗi ngày tổ làm $36$ sản phẩm nên hoàn thành sớm $2$ ngày và còn làm vượt $12$ sản phẩm. Hỏi theo kế hoạch tổ phải làm bao nhiêu sản phẩm?
\end{probbox}
\probsol{
Nếu làm đủ số ngày dự định với $36$ sản phẩm mỗi ngày thì làm được nhiều hơn kế hoạch $36\times2+12=84$ sản phẩm.\par
Mỗi ngày làm vượt $6$ sản phẩm, nên số ngày dự định là $84:6=14$ ngày, tức là\par $30\times14=420$ sản phẩm.
\dapso{$420$ sản phẩm.}}
\gap
\begin{probbox}{34}{\stars{1}}
An mua $5$ bút xanh và $3$ bút đỏ hết $33\,000$ đồng. Bình mua $3$ bút xanh và $5$ bút đỏ cùng loại hết $31\,000$ đồng. Tìm giá mỗi loại bút.
\end{probbox}
\probsol{
Cộng lại: $8$ bút xanh và $8$ bút đỏ giá $64\,000$ đồng, nên $1$ bút xanh $+$ $1$ bút đỏ giá $8\,000$ đồng.\par
So với Bình, An mua nhiều hơn $2$ bút xanh, ít hơn $2$ bút đỏ và trả nhiều hơn $2\,000$ đồng, nên một bút xanh đắt hơn một bút đỏ $1\,000$ đồng.\par
Bút xanh: $(8\,000+1\,000):2=4\,500$ đồng; bút đỏ: $3\,500$ đồng.
\dapso{Bút xanh: $4\,500$ đồng; bút đỏ: $3\,500$ đồng.}}
\gap
\begin{probbox}{35}{\stars{2}}
Một người được thuê làm trong $30$ ngày. Mỗi ngày làm việc được trả $200\,000$ đồng, mỗi ngày nghỉ bị trừ $50\,000$ đồng. Sau $30$ ngày người đó nhận được $4\,750\,000$ đồng. Hỏi người đó đã nghỉ bao nhiêu ngày?
\end{probbox}
\probsol{
Nếu làm đủ $30$ ngày thì được $6\,000\,000$ đồng, nhiều hơn số tiền thực nhận $1\,250\,000$ đồng.\par
Mỗi ngày nghỉ làm giảm $200\,000+50\,000=250\,000$ đồng, nên số ngày nghỉ là $1\,250\,000:250\,000=5$.
\dapso{$5$ ngày.}}
\gap
\begin{probbox}{36}{\stars{2}}
Một cuộc thi có $20$ bài. Mỗi bài làm đúng được $5$ điểm; mỗi bài làm sai hoặc bỏ qua bị trừ $2$ điểm. Minh được $72$ điểm. Hỏi Minh làm đúng bao nhiêu bài?
\end{probbox}
\probsol{
Nếu đúng cả $20$ bài thì được $100$ điểm, nhiều hơn $28$ điểm.\par Mỗi bài sai làm mất $5+2=7$ điểm, nên có $28:7=4$ bài sai và $16$ bài đúng.
\dapso{$16$ bài.}}
\gap
\begin{probbox}{37}{\stars{1}}
Chia kẹo: mỗi em $5$ cái thì thừa $7$ cái; mỗi em $6$ cái thì có một em chỉ được $2$ cái. Tìm số em và số kẹo.
\end{probbox}
\probsol{
Khi chia mỗi em $6$ cái, em cuối cùng thiếu $6-2=4$ cái.\par Hai cách chia chênh nhau $7+4=11$; mỗi em chênh nhau $1$.\par Số em: $11$.\par Số kẹo: $5\times11+7=62$.
\dapso{$11$ em, $62$ cái kẹo.}}
\gap
\begin{probbox}{38}{\stars{1}}
Có $38$ con gà và chó. Số chân gà nhiều hơn số chân chó là $16$ chân. Hỏi mỗi loại có bao nhiêu con?
\end{probbox}
\probsol{
Nếu cả $38$ con đều là gà thì hiệu số chân là $2\times38=76$ chân.\par Như vậy nhiều hơn thực tế $76-16=60$.\par Thay một con gà bằng một con chó thì hiệu giảm $2+4=6$.\par Số chó: $60:6=10$; số gà: $38-10=28$.\par Thử lại: $56-40=16$.
\dapso{$28$ con gà, $10$ con chó.}}
\gap
\begin{probbox}{39}{\stars{1}}
Một người được thuê làm trong $26$ ngày: mỗi ngày làm việc được trả $300\,000$ đồng, mỗi ngày nghỉ bị phạt $100\,000$ đồng. Người đó nhận được $6\,200\,000$ đồng. Hỏi người đó đã nghỉ bao nhiêu ngày?
\end{probbox}
\probsol{
Nếu làm đủ $26$ ngày: $300\,000\times26=7\,800\,000$ đồng.\par Số tiền bị hụt: $7\,800\,000-6\,200\,000=1\,600\,000$ đồng.\par Mỗi ngày nghỉ làm mất $300\,000+100\,000=400\,000$ đồng.\par Số ngày nghỉ: $1\,600\,000:400\,000=4$.
\dapso{$4$ ngày.}}
\gap
\begin{probbox}{40}{\stars{1}}
Một ô tô đi từ A. Nếu đi với vận tốc $40$ km/giờ thì đến B lúc $11$ giờ; nếu đi với vận tốc $50$ km/giờ thì đến B lúc $10$ giờ. Tìm quãng đường AB và thời điểm xuất phát.
\end{probbox}
\probsol{
Đi với vận tốc $40$ km/giờ thì mất nhiều hơn $11-10=1$ giờ.\par Nếu xe đi nhanh tiếp tục đi thêm giờ đó sau khi đến B thì sẽ vượt quá B $50\times1=50$ km.\par Vậy trong cùng thời gian, xe đi nhanh đi được nhiều hơn $50$ km, mỗi giờ hơn $50-40=10$ km.\par Thời gian đi với vận tốc $40$ km/giờ: $50:10=5$ giờ.\par AB: $40\times5=200$ km.\par Xuất phát lúc: $11-5=6$ giờ.
\dapso{$200$ km; xuất phát lúc $6$ giờ.}}
\gap
\begin{probbox}{41}{\stars{1}}
Học sinh ngồi vào bàn. Nếu mỗi bàn $2$ em thì $8$ em không có chỗ; nếu mỗi bàn $3$ em thì thừa $1$ bàn và có một bàn chỉ có $1$ em. Tìm số bàn và số học sinh.
\end{probbox}
\probsol{
Khi xếp mỗi bàn $3$ em, bàn trống thiếu $3$ em và bàn kia thiếu $2$ em: tổng cộng thiếu $5$ em.\par Hai cách xếp chênh nhau $8+5=13$ học sinh; mỗi bàn chênh nhau $1$.\par Số bàn: $13$.\par Số học sinh: $2\times13+8=34$.
\dapso{$13$ bàn, $34$ học sinh.}}
\gap
\begin{probbox}{42}{\stars{1}}
Một bài kiểm tra có $25$ câu: mỗi câu đúng được $4$ điểm, mỗi câu sai bị trừ $1$ điểm. Hoa trả lời hết các câu và được $70$ điểm. Hỏi Hoa trả lời sai bao nhiêu câu?
\end{probbox}
\probsol{
Nếu đúng tất cả thì được $4\times25=100$ điểm.\par Số điểm bị hụt: $100-70=30$.\par Mỗi câu sai làm mất $4+1=5$ điểm.\par Số câu sai: $30:5=6$.
\dapso{$6$ câu.}}
\gap
\begin{probbox}{43}{\stars{1}}
Mua $2$ áo sơ mi và $3$ quần dài hết $1\,100\,000$ đồng; mua $3$ áo sơ mi và $2$ quần dài hết $1\,150\,000$ đồng. Tìm giá mỗi loại.
\end{probbox}
\probsol{
Cộng lại: $5$ áo và $5$ quần giá $2\,250\,000$ đồng.\par Một áo và một quần: $2\,250\,000:5=450\,000$ đồng.\par Lần thứ hai mua nhiều hơn một áo, ít hơn một quần và trả nhiều hơn $50\,000$ đồng, nên một áo đắt hơn một quần $50\,000$ đồng.\par Áo: $(450\,000+50\,000):2=250\,000$ đồng.\par Quần: $450\,000-250\,000=200\,000$ đồng.
\dapso{Áo $250\,000$ đồng, quần $200\,000$ đồng.}}
\gap
\begin{probbox}{44}{\stars{1}}
Mẹ mua một số cam. Nếu bán với giá $30\,000$ đồng/kg thì lãi $40\,000$ đồng; nếu bán với giá $25\,000$ đồng/kg thì lỗ $10\,000$ đồng. Tìm số cam và số tiền vốn.
\end{probbox}
\probsol{
Hai mức giá cho số tiền chênh nhau $40\,000+10\,000=50\,000$ đồng.\par Mỗi ki-lô-gam chênh: $30\,000-25\,000=5\,000$ đồng.\par Số cam: $50\,000:5\,000=10$ kg.\par Tiền vốn: $30\,000\times10-40\,000=260\,000$ đồng.
\dapso{$10$ kg; tiền vốn $260\,000$ đồng.}}
\gap
\begin{probbox}{45}{\stars{1}}
$18$ con nhện ($8$ chân, không có cánh) và chuồn chuồn ($6$ chân, $2$ đôi cánh) có $124$ chân. Hỏi có bao nhiêu đôi cánh?
\end{probbox}
\probsol{
Nếu tất cả đều là chuồn chuồn: $6\times18=108$ chân.\par Còn thiếu $124-108=16$; mỗi con nhện nhiều hơn $2$ chân.\par Số nhện: $8$; số chuồn chuồn: $10$.\par Số đôi cánh: $2\times10=20$.
\dapso{$20$ đôi cánh.}}
\gap
\begin{probbox}{46}{\stars{1}}
Một cửa hàng có $45$ hộp bánh, mỗi hộp đựng $6$ hoặc $10$ cái, tất cả có $350$ cái bánh. Hỏi mỗi loại có bao nhiêu hộp?
\end{probbox}
\probsol{
Nếu tất cả các hộp đều đựng $6$ cái: $6\times45=270$.\par Còn thiếu $80$; mỗi hộp $10$ cái nhiều hơn $4$ cái.\par Số hộp $10$ cái: $80:4=20$.\par Số hộp $6$ cái: $25$.
\dapso{$25$ hộp loại $6$ cái, $20$ hộp loại $10$ cái.}}
\gap
\begin{probbox}{47}{\stars{2}}
An đi xe đạp đến trường. Nếu đi với vận tốc $12$ km/giờ thì đến muộn $5$ phút; nếu đi với vận tốc $15$ km/giờ thì đến sớm $5$ phút. Hỏi trường cách nhà bao xa?
\end{probbox}
\probsol{
$5$ phút $=\dfrac1{12}$ giờ.\par Đi $12$ km/giờ trong thời gian quy định thì còn thiếu $12\times\dfrac1{12}=1$ km.\par Đi $15$ km/giờ thì vượt quá $15\times\dfrac1{12}=1{,}25$ km.\par Hiệu số: $2{,}25$ km, mỗi giờ chênh $3$ km.\par Thời gian quy định: $2{,}25:3=0{,}75$ giờ.\par Quãng đường: $12\times0{,}75+1=10$ km.
\dapso{$10$ km.}}
\gap
\begin{probbox}{48}{\stars{1}}
Một bãi đỗ xe có $50$ chiếc xe gồm ô tô ($4$ bánh) và xe máy ($2$ bánh). Số bánh ô tô nhiều hơn số bánh xe máy là $20$ bánh. Hỏi mỗi loại có bao nhiêu chiếc?
\end{probbox}
\probsol{
Nếu cả $50$ chiếc đều là ô tô thì hiệu là $4\times50=200$.\par Như vậy nhiều hơn thực tế $180$.\par Thay một ô tô bằng một xe máy thì hiệu giảm $4+2=6$.\par Số xe máy: $180:6=30$; số ô tô: $20$.\par Thử lại: $80-60=20$.
\dapso{$20$ ô tô, $30$ xe máy.}}
\gap
\begin{probbox}{49}{\stars{2}}
Một bài kiểm tra có $30$ câu: câu đúng được $5$ điểm, câu sai bị trừ $2$ điểm, câu bỏ qua được $0$ điểm. Bình được $112$ điểm, và số câu sai gấp đôi số câu bỏ qua. Hỏi Bình trả lời đúng bao nhiêu câu?
\end{probbox}
\probsol{
Ghép $1$ câu bỏ qua với $2$ câu sai thành một nhóm $3$ câu.\par Nếu đúng tất cả: $5\times30=150$; số điểm bị hụt: $150-112=38$.\par Mỗi nhóm làm mất $5+(5+2)\times2=19$ điểm.\par Số nhóm: $38:19=2$, tức là\par $2$ câu bỏ qua và $4$ câu sai.\par Số câu đúng: $30-2-4=24$.
\dapso{$24$ câu.}}
\gap
\begin{probbox}{50}{\stars{1}}
Hùng mang một số tiền đi mua vở. Nếu mua $8$ quyển thì còn thừa $4\,000$ đồng; nếu mua $10$ quyển thì thiếu $6\,000$ đồng. Hỏi Hùng mang bao nhiêu tiền?
\end{probbox}
\probsol{
Hai cách mua chênh nhau $4\,000+6\,000=10\,000$ đồng ứng với $2$ quyển vở.\par Một quyển vở: $5\,000$ đồng.\par Số tiền: $5\,000\times8+4\,000=44\,000$ đồng.
\dapso{$44\,000$ đồng.}}
\clearpage
\titlehead
\begingroup\setbox0=\vbox{\section*{Tính ngược từ cuối}}\endgroup
\chapterhead{10}{Tính ngược từ cuối}

\shead{Bắt đầu từ kết quả cuối cùng}
\begin{knowbox}
Khi bài toán cho \textbf{kết quả cuối cùng} của một chuỗi phép tính (hoặc một chuỗi thay đổi), ta bắt đầu từ kết quả đó và \textbf{tính ngược lại} từng bước, mỗi bước dùng \emph{phép tính ngược}:
\begin{center}
cộng \ $\longleftrightarrow$ \ trừ \qquad\qquad nhân \ $\longleftrightarrow$ \ chia
\end{center}
\begin{itemize}[leftmargin=1.4em,itemsep=3pt]
\item Vẽ \textbf{sơ đồ} các bước để không bỏ sót hoặc làm sai thứ tự bước nào.
\item Với các bài ``lấy đi một phân số, rồi thêm (hoặc bớt) vài đơn vị nữa'', ở mỗi bước tính ngược, ta xử lí \emph{số đơn vị thêm} trước rồi mới đến \emph{phân số}.
\item Với các bài có số lượng chuyển qua lại giữa nhiều người (hoặc nhiều bình), lập \textbf{bảng} ghi số lượng sau mỗi lần chuyển và tính từ hàng dưới cùng lên.
\end{itemize}
Luôn \textbf{thử lại} bằng cách tính xuôi.
\end{knowbox}
\gap
\begin{exbox}{\ExWord\ 1}
Tìm một số, biết rằng lấy số đó cộng với $15$, nhân với $3$, trừ đi $12$ rồi cuối cùng chia cho $4$ thì được $18$.
\guide
Sơ đồ (tính xuôi ở trên, tính ngược ở dưới):
\begin{center}
\begin{tikzpicture}[x=2.25cm,box/.style={draw=tmBlue!70!black,line width=0.8pt,rounded corners=2pt,minimum width=1.05cm,minimum height=0.62cm,font=\small}]
\foreach \v [count=\i from 0] in {?,28+15,84,72,18}{\node[box] (n\i) at (\i,0) {};}
\node at (n0) {$\mathbf{13}$}; \node at (n1) {$28$}; \node at (n2) {$84$}; \node at (n3) {$72$}; \node[box,fill=tmOrange!25] at (n4) {$18$};
\foreach \a/\b/\f/\g in {0/1/+15/-15,1/2/\times3/:3,2/3/-12/+12,3/4/:4/\times4}{
  \draw[-{Stealth},line width=0.7pt] ([yshift=3pt]n\a.east) to[bend left=28] node[above,font=\footnotesize]{$\f$} ([yshift=3pt]n\b.west);
  \draw[-{Stealth},line width=0.7pt,tmOrange!80!black] ([yshift=-3pt]n\b.west) to[bend left=28] node[below,font=\footnotesize]{$\g$} ([yshift=-3pt]n\a.east);}
\end{tikzpicture}
\end{center}
Tính ngược: $18\times4=72$;\par $72+12=84$;\par $84:3=28$;\par $28-15=13$.\par
\emph{Thử lại:} $(13+15)\times3=84$;\par $84-12=72$;\par $72:4=18$.
\dapso{$13$.}
\end{exbox}
\gap
\begin{exbox}{\ExWord\ 2}
Một bà mang trứng ra chợ bán. Lần đầu bà bán một nửa số trứng và thêm $1$ quả. Lần sau bà bán một nửa số trứng còn lại và thêm $1$ quả. Bà còn lại $10$ quả. Hỏi bà mang ra chợ bao nhiêu quả trứng?
\guide
\emph{Tính ngược lần bán thứ hai:} nếu không bán thêm quả trứng đó thì còn $10+1=11$ quả, đó là một nửa số trứng trước lần bán này.\par Vậy trước lần bán thứ hai: $11\times2=22$ quả.\par
\emph{Tính ngược lần bán thứ nhất:} $(22+1)\times2=46$ quả.
\dapso{$46$ quả trứng.}
\end{exbox}
\gap
\begin{exbox}{\ExWord\ 3}
Ngày thứ nhất dùng $\dfrac13$ lượng nước trong bể; ngày thứ hai dùng $\dfrac25$ lượng nước còn lại. Bây giờ bể còn $120$ lít. Hỏi lúc đầu bể có bao nhiêu lít nước?
\guide
Sau ngày thứ hai, còn lại $1-\dfrac25=\dfrac35$ lượng nước còn sau ngày thứ nhất.\par Lượng nước còn sau ngày thứ nhất là $120:\dfrac35=200$ lít.\par
$200$ lít là $1-\dfrac13=\dfrac23$ lượng nước ban đầu.\par Lúc đầu: $200:\dfrac23=300$ lít.
\dapso{$300$ lít.}
\end{exbox}
\gap
\begin{exbox}{\ExWord\ 4}
Ba thùng chứa $240$ lít dầu. Đầu tiên, thùng I rót sang thùng II một lượng dầu bằng lượng dầu thùng II đang có. Sau đó thùng II rót sang thùng III một lượng bằng lượng dầu thùng III đang có. Rồi thùng III rót sang thùng I một lượng bằng lượng dầu thùng I đang có. Lúc này ba thùng có lượng dầu bằng nhau. Hỏi lúc đầu mỗi thùng có bao nhiêu lít dầu?
\guide
Lúc cuối mỗi thùng có $240:3=80$ lít.\par Thùng nào ``nhận thêm bằng lượng đang có'' thì gấp đôi; khi tính ngược, ta \emph{chia đôi} thùng nhận và \emph{trả lại} lượng đó cho thùng rót.
\begin{center}
\renewcommand{\arraystretch}{1.25}
\begin{tabular}{|l|c|c|c|}
\hline
 & Thùng I & Thùng II & Thùng III\\\hline
Lúc cuối & $80$ & $80$ & $80$\\\hline
Trước bước 3 & $40$ & $80$ & $120$\\\hline
Trước bước 2 & $40$ & $140$ & $60$\\\hline
Lúc đầu & $\mathbf{110}$ & $\mathbf{70}$ & $\mathbf{60}$\\\hline
\end{tabular}
\end{center}
\dapso{Thùng I: $110$ lít; thùng II: $70$ lít; thùng III: $60$ lít.}
\end{exbox}
\dividerbar{Luyện tập}
\begin{probbox}{1}{\relax}
Tìm một số, biết rằng lấy số đó nhân với $4$, trừ đi $18$ rồi chia cho $2$ thì được $25$.
\end{probbox}
\probsol{
$25\times2=50$;\par $50+18=68$;\par $68:4=17$.
\dapso{$17$.}}
\gap
\begin{probbox}{2}{\relax}
Nam nghĩ ra một số. Bạn ấy nhân số đó với $6$, cộng thêm $8$, chia cho $2$ rồi trừ đi $14$ thì được $20$. Hỏi Nam đã nghĩ ra số nào?
\end{probbox}
\probsol{
$20+14=34$;\par $34\times2=68$;\par $68-8=60$;\par $60:6=10$.
\dapso{$10$.}}
\gap
\begin{probbox}{3}{\relax}
Tìm một số, biết rằng lấy số đó cộng với $\dfrac12$ rồi nhân với $\dfrac23$ thì được $3$.
\end{probbox}
\probsol{
$3:\dfrac23=\dfrac92$;\par $\dfrac92-\dfrac12=4$.
\dapso{$4$.}}
\gap
\begin{probbox}{4}{\relax}
Lấy tuổi của An gấp đôi, cộng thêm $5$ rồi chia cho $3$ thì được $9$. Hỏi An bao nhiêu tuổi?
\end{probbox}
\probsol{
$9\times3=27$;\par $27-5=22$;\par $22:2=11$.
\dapso{$11$ tuổi.}}
\gap
\begin{probbox}{5}{\relax}
Lan dùng $\dfrac12$ số tiền để mua sách và $\dfrac13$ số tiền còn lại để mua bút, còn lại $20\,000$ đồng. Hỏi lúc đầu Lan có bao nhiêu tiền?
\end{probbox}
\probsol{
$20\,000$ đồng là $\dfrac23$ số tiền còn lại sau khi mua sách, tức là\par $30\,000$ đồng.\par Lúc đầu: $30\,000\times2=60\,000$ đồng.
\dapso{$60\,000$ đồng.}}
\gap
\begin{probbox}{6}{\relax}
Ngày thứ nhất Hoa đọc $\dfrac14$ cuốn sách, ngày thứ hai đọc $\dfrac13$ số trang còn lại, ngày thứ ba đọc nốt $60$ trang cuối. Hỏi cuốn sách có bao nhiêu trang?
\end{probbox}
\probsol{
$60$ trang là $\dfrac23$ số trang còn lại sau ngày thứ nhất: $60:\dfrac23=90$.\par Và $90$ trang là $\dfrac34$ cuốn sách: $90:\dfrac34=120$.
\dapso{$120$ trang.}}
\gap
\begin{probbox}{7}{\relax}
Bác Ba bán một nửa số gà và thêm $2$ con, rồi bán một nửa số gà còn lại và thêm $2$ con, thì còn lại $8$ con. Hỏi lúc đầu bác Ba có bao nhiêu con gà?
\end{probbox}
\probsol{
Trước lần bán thứ hai: $(8+2)\times2=20$.\par Lúc đầu: $(20+2)\times2=44$.
\dapso{$44$ con gà.}}
\gap
\begin{probbox}{8}{\relax}
Đi chợ, lần đầu mẹ tiêu một nửa số tiền và thêm $5\,000$ đồng, lần sau tiêu một nửa số tiền còn lại và thêm $5\,000$ đồng, thì còn lại $10\,000$ đồng. Hỏi mẹ mang đi bao nhiêu tiền?
\end{probbox}
\probsol{
Trước lần mua thứ hai: $(10\,000+5\,000)\times2=30\,000$ đồng.\par Lúc đầu: $(30\,000+5\,000)\times2=70\,000$ đồng.
\dapso{$70\,000$ đồng.}}
\gap
\begin{probbox}{9}{\relax}
Rót $15$ lít từ bể A sang bể B, rồi rót $5$ lít từ bể B trở lại bể A. Lúc này mỗi bể có $40$ lít. Hỏi lúc đầu mỗi bể có bao nhiêu lít?
\end{probbox}
\probsol{
Bể A: $40-5+15=50$ lít; bể B: $40+5-15=30$ lít.
\dapso{A: $50$ lít; B: $30$ lít.}}
\gap
\begin{probbox}{10}{\relax}
Chuyển $6$ quả cam từ rổ I sang rổ II, rồi chuyển $2$ quả từ rổ II trở lại rổ I. Lúc này rổ I có $20$ quả và rổ II có $30$ quả. Hỏi lúc đầu mỗi rổ có bao nhiêu quả?
\end{probbox}
\probsol{
Rổ I: $20-2+6=24$; rổ II: $30+2-6=26$.
\dapso{$24$ quả và $26$ quả cam.}}
\gap
\begin{probbox}{11}{\relax}
Bình I rót một nửa số nước sang bình II; sau đó bình II rót $\dfrac13$ lượng nước đang có sang bình I. Lúc này bình I có $8$ lít và bình II có $12$ lít. Hỏi lúc đầu mỗi bình có bao nhiêu lít nước?
\end{probbox}
\probsol{
Trước lần rót thứ hai, bình II còn giữ $\dfrac23$, nên lúc đó có $12:\dfrac23=18$ lít và đã rót đi $6$ lít; bình I có $8-6=2$ lít.\par
Lúc đầu bình I có $2\times2=4$ lít và đã rót đi $2$ lít; bình II có $18-2=16$ lít.
\dapso{Bình I: $4$ lít; bình II: $16$ lít.}}
\gap
\begin{probbox}{12}{\relax}
Tìm một số: cộng với $7$, nhân với $5$, rồi trừ đi $10$ thì được $65$.
\end{probbox}
\probsol{
Trước khi trừ $10$: $65+10=75$.\par Trước khi nhân với $5$: $75:5=15$.\par Trước khi cộng $7$: $15-7=8$.
\dapso{$8$.}}
\gap
\begin{probbox}{13}{\relax}
Tìm một số: nhân với $3$, cộng thêm $12$, rồi chia cho $5$ thì được $9$.
\end{probbox}
\probsol{
Trước khi chia cho $5$: $9\times5=45$.\par Trước khi cộng $12$: $45-12=33$.\par Số cần tìm: $33:3=11$.
\dapso{$11$.}}
\gap
\begin{probbox}{14}{\relax}
Tìm một số: trừ đi $4$, nhân với $6$, rồi cộng thêm $3$ thì được $51$.
\end{probbox}
\probsol{
Trước khi cộng $3$: $51-3=48$.\par Trước khi nhân với $6$: $48:6=8$.\par Số cần tìm: $8+4=12$.
\dapso{$12$.}}
\gap
\begin{probbox}{15}{\relax}
Tìm một số: chia cho $2$, cộng thêm $15$, rồi nhân với $4$ thì được $100$.
\end{probbox}
\probsol{
Trước khi nhân với $4$: $100:4=25$.\par Trước khi cộng $15$: $25-15=10$.\par Số cần tìm: $10\times2=20$.
\dapso{$20$.}}
\gap
\begin{probbox}{16}{\relax}
Tìm $x$: nhân với $2{,}5$, rồi trừ đi $1{,}5$ thì được $8{,}5$.
\end{probbox}
\probsol{
Trước khi trừ $1{,}5$: $8{,}5+1{,}5=10$.\par $x=10:2{,}5=4$.
\dapso{$x=4$.}}
\gap
\begin{probbox}{17}{\relax}
Tìm một số thập phân: cộng thêm $2{,}4$, rồi nhân với $3$ thì được $12{,}6$.
\end{probbox}
\probsol{
Trước khi nhân với $3$: $12{,}6:3=4{,}2$.\par Số cần tìm: $4{,}2-2{,}4=1{,}8$.
\dapso{$1{,}8$.}}
\gap
\begin{probbox}{18}{\relax}
Mai mua một cuốn sách hết $25\,000$ đồng, được mẹ cho thêm $10\,000$ đồng, rồi mua một cái bút hết $8\,000$ đồng thì còn lại $17\,000$ đồng. Hỏi lúc đầu Mai có bao nhiêu tiền?
\end{probbox}
\probsol{
Trước khi mua bút: $17\,000+8\,000=25\,000$ đồng.\par Trước khi được mẹ cho: $25\,000-10\,000=15\,000$ đồng.\par Lúc đầu: $15\,000+25\,000=40\,000$ đồng.
\dapso{$40\,000$ đồng.}}
\gap
\begin{probbox}{19}{\relax}
Một xe buýt rời bến. Đến điểm dừng thứ nhất có $6$ người xuống và $4$ người lên; đến điểm dừng thứ hai có $10$ người xuống và $3$ người lên. Lúc này trên xe có $21$ hành khách. Hỏi lúc rời bến trên xe có bao nhiêu hành khách?
\end{probbox}
\probsol{
Trước điểm dừng thứ hai: $21-3+10=28$.\par Trước điểm dừng thứ nhất: $28-4+6=30$.
\dapso{$30$ hành khách.}}
\gap
\begin{probbox}{20}{\relax}
Lấy tuổi bố chia cho $4$ rồi cộng thêm $3$ thì được $13$. Hỏi bố bao nhiêu tuổi?
\end{probbox}
\probsol{
Trước khi cộng $3$: $13-3=10$.\par Tuổi bố: $10\times4=40$.
\dapso{$40$ tuổi.}}
\gap
\begin{probbox}{21}{\relax}
Lan tiêu $\dfrac13$ số tiền, rồi tiêu $\dfrac12$ số tiền còn lại, thì còn $25\,000$ đồng. Hỏi lúc đầu Lan có bao nhiêu tiền?
\end{probbox}
\probsol{
$25\,000$ đồng là $\dfrac12$ số tiền còn lại sau lần tiêu thứ nhất.\par Vậy sau lần tiêu thứ nhất còn $25\,000\times2=50\,000$ đồng.\par Số tiền đó là $\dfrac23$ số tiền ban đầu.\par Ban đầu: $50\,000:\dfrac23=75\,000$ đồng.
\dapso{$75\,000$ đồng.}}
\gap
\begin{probbox}{22}{\relax}
Ngày thứ nhất dùng $\dfrac14$ lượng nước trong bể; ngày thứ hai dùng $\dfrac13$ lượng nước còn lại. Bể còn $300$ lít. Hỏi lúc đầu bể có bao nhiêu lít nước?
\end{probbox}
\probsol{
$300$ lít là $\dfrac23$ lượng nước còn lại sau ngày thứ nhất.\par Sau ngày thứ nhất: $300:\dfrac23=450$ lít.\par Đó là $\dfrac34$ lượng nước ban đầu.\par Ban đầu: $450:\dfrac34=600$ lít.
\dapso{$600$ lít.}}
\gap
\begin{probbox}{23}{\relax}
Một người bán trứng, lần đầu bán một nửa số trứng và thêm $2$ quả, lần sau bán một nửa số còn lại và thêm $1$ quả, thì còn lại $9$ quả. Hỏi lúc đầu có bao nhiêu quả trứng?
\end{probbox}
\probsol{
Trước khi bán thêm $1$ quả: $9+1=10$, là một nửa số trứng còn lại sau lần bán thứ nhất.\par Sau lần bán thứ nhất: $10\times2=20$.\par Tương tự, $20+2=22$ là một nửa số trứng ban đầu.\par Ban đầu: $22\times2=44$.
\dapso{$44$ quả trứng.}}
\gap
\begin{probbox}{24}{\relax}
Chuyển $8$ quả bưởi từ rổ A sang rổ B, rồi chuyển $3$ quả từ rổ B sang rổ A. Lúc này rổ A có $25$ quả và rổ B có $35$ quả. Hỏi lúc đầu mỗi rổ có bao nhiêu quả?
\end{probbox}
\probsol{
Trước lần chuyển thứ hai: A $25-3=22$, B $35+3=38$.\par Lúc đầu: A $22+8=30$, B $38-8=30$.
\dapso{Mỗi rổ $30$ quả.}}
\gap
\begin{probbox}{25}{\relax}
Một cửa hàng ngày thứ nhất bán $\dfrac25$ số gạo, ngày thứ hai bán một nửa số gạo còn lại, thì còn $45$ kg. Hỏi lúc đầu cửa hàng có bao nhiêu ki-lô-gam gạo?
\end{probbox}
\probsol{
Sau ngày thứ nhất: $45\times2=90$ kg.\par Đó là $\dfrac35$ số gạo ban đầu.\par Ban đầu: $90:\dfrac35=150$ kg.
\dapso{$150$ kg.}}
\gap
\begin{probbox}{26}{\relax}
Tìm $x$, biết $\left(x-\dfrac13\right)\times\dfrac34=\dfrac12$.
\end{probbox}
\probsol{
$x-\dfrac13=\dfrac12:\dfrac34=\dfrac23$.\par $x=\dfrac23+\dfrac13=1$.
\dapso{$x=1$.}}
\gap
\begin{probbox}{27}{\relax}
An cho Bình $\dfrac13$ số bi của mình, rồi cho Chi $4$ viên thì còn lại $12$ viên. Hỏi lúc đầu An có bao nhiêu viên bi?
\end{probbox}
\probsol{
Trước khi cho Chi: $12+4=16$.\par Đó là $\dfrac23$ số bi ban đầu.\par Ban đầu: $16:\dfrac23=24$.
\dapso{$24$ viên bi.}}
\gap
\begin{probbox}{28}{\relax}
Tìm một số: nhân với $4$, cộng thêm $0{,}6$, rồi chia cho $2$ thì được $5{,}3$.
\end{probbox}
\probsol{
Trước khi chia cho $2$: $5{,}3\times2=10{,}6$.\par Trước khi cộng $0{,}6$: $10$.\par Số cần tìm: $10:4=2{,}5$.
\dapso{$2{,}5$.}}
\gap
\begin{probbox}{29}{\relax}
Từ một cuộn vải, lần đầu cắt đi $\dfrac14$, lần thứ hai cắt $5$ m, lần thứ ba cắt một nửa phần còn lại, thì còn $8$ m. Hỏi cuộn vải dài bao nhiêu mét?
\end{probbox}
\probsol{
Trước lần cắt thứ ba: $8\times2=16$ m.\par Trước lần cắt thứ hai: $16+5=21$ m.\par Đó là $\dfrac34$ cuộn vải.\par Cuộn vải: $21:\dfrac34=28$ m.
\dapso{$28$ m.}}
\gap
\begin{probbox}{30}{\relax}
Chú Tư bán $\dfrac13$ số vịt, rồi mua thêm $15$ con thì có $75$ con. Hỏi lúc đầu chú Tư có bao nhiêu con vịt?
\end{probbox}
\probsol{
Trước khi mua: $75-15=60$.\par Đó là $\dfrac23$ đàn vịt.\par Ban đầu: $60:\dfrac23=90$.
\dapso{$90$ con vịt.}}
\dividerbar{Thử thách -- Nâng cao}
\begin{probbox}{31}{\stars{1}}
Một cửa hàng bán $\dfrac13$ cuộn vải và thêm $2$ m, rồi bán một nửa phần còn lại và thêm $3$ m, thì còn $12$ m. Hỏi cuộn vải dài bao nhiêu mét?
\end{probbox}
\probsol{
Trước lần bán thứ hai: $(12+3)\times2=30$ m.\par Khi đó $30+2=32$ m là $\dfrac23$ cuộn vải, nên cuộn vải dài $32:\dfrac23=48$ m.
\dapso{$48$ m.}}
\gap
\begin{probbox}{32}{\stars{1}}
Một chú khỉ hái được một số quả đào. Mỗi ngày khỉ ăn một nửa số đào đang có và thêm $1$ quả. Sáng ngày thứ năm khỉ chỉ còn đúng $1$ quả. Hỏi khỉ đã hái được bao nhiêu quả đào?
\end{probbox}
\probsol{
Khỉ đã ăn trong $4$ ngày.\par Tính ngược mỗi ngày bằng (số còn lại $+1$)$\times2$.\par
$1\to4\to10\to22\to46$.
\dapso{$46$ quả đào.}}
\gap
\begin{probbox}{33}{\stars{1}}
An, Bình và Cường có $96$ viên bi. An cho Bình số bi bằng số bi Bình đang có; rồi Bình cho Cường số bi bằng số bi Cường đang có; cuối cùng Cường cho An số bi bằng số bi An đang có. Lúc này số bi của ba bạn bằng nhau. Hỏi lúc đầu mỗi bạn có bao nhiêu viên bi?
\end{probbox}
\probsol{
Lúc cuối mỗi bạn có $32$ viên.\par Tính ngược từ cuối (như Ví dụ $4$):\par
Trước bước 3: An $16$, Bình $32$, Cường $48$.\par Trước bước 2: An $16$, Bình $56$, Cường $24$.\par Lúc đầu: An $44$, Bình $28$, Cường $24$.
\dapso{An: $44$; Bình: $28$; Cường: $24$ viên bi.}}
\gap
\begin{probbox}{34}{\stars{2}}
A, B và C chơi ba ván. Sau mỗi ván, người thua phải cho mỗi người còn lại số bi bằng số bi người đó đang có. A thua ván 1, B thua ván 2 và C thua ván 3. Sau đó mỗi người có $24$ viên bi. Hỏi lúc đầu mỗi người có bao nhiêu viên bi?
\end{probbox}
\probsol{
Tổng cộng: $72$ viên bi.\par Tính ngược, ở mỗi ván hai người thắng chia đôi số bi và người thua nhận lại số bi đó.\par
Trước ván 3: A $12$, B $12$, C $48$.\par Trước ván 2: A $6$, B $42$, C $24$.\par Lúc đầu: A $39$, B $21$, C $12$.
\dapso{A: $39$; B: $21$; C: $12$ viên bi.}}
\gap
\begin{probbox}{35}{\stars{1}}
Mỗi ngày một chú khỉ ăn một nửa số đào đang có và thêm $2$ quả. Sau $3$ ngày khỉ còn $3$ quả. Hỏi lúc đầu khỉ có bao nhiêu quả đào?
\end{probbox}
\probsol{
Đầu ngày thứ ba: $(3+2)\times2=10$.\par Đầu ngày thứ hai: $(10+2)\times2=24$.\par Lúc đầu: $(24+2)\times2=52$.
\dapso{$52$ quả đào.}}
\gap
\begin{probbox}{36}{\stars{1}}
Rót từ bình A sang bình B một lượng bằng lượng nước bình B đang có; rồi rót từ bình B sang bình A một lượng bằng lượng nước bình A đang có. Lúc này mỗi bình có $24$ lít. Hỏi lúc đầu mỗi bình có bao nhiêu lít?
\end{probbox}
\probsol{
Lần rót thứ hai làm bình A gấp đôi, nên trước đó A có $24:2=12$ và B có $24+12=36$.\par Lần rót thứ nhất làm bình B gấp đôi, nên trước đó B có $36:2=18$ và A có $12+18=30$.
\dapso{A $30$ lít, B $18$ lít.}}
\gap
\begin{probbox}{37}{\stars{1}}
Một cửa hàng bán một nửa số gạo và thêm $5$ kg, rồi bán $\dfrac13$ số gạo còn lại và thêm $4$ kg, rồi bán $\dfrac14$ số gạo còn lại và thêm $3$ kg, thì còn $15$ kg. Hỏi lúc đầu có bao nhiêu ki-lô-gam gạo?
\end{probbox}
\probsol{
Trước lần bán thứ ba: $15+3=18$ kg là $\dfrac34$, nên có $18:\dfrac34=24$ kg.\par Trước lần bán thứ hai: $24+4=28$ kg là $\dfrac23$, nên có $42$ kg.\par Lúc đầu: $(42+5)\times2=94$ kg.
\dapso{$94$ kg.}}
\gap
\begin{probbox}{38}{\stars{1}}
Nam nghĩ ra một số. Ba lần liên tiếp, bạn ấy gấp đôi số đó rồi trừ đi $3$, cuối cùng được $19$. Hỏi Nam đã nghĩ ra số nào?
\end{probbox}
\probsol{
Trước lượt thứ ba: $(19+3):2=11$.\par Trước lượt thứ hai: $(11+3):2=7$.\par Số cần tìm: $(7+3):2=5$.
\dapso{$5$.}}
\gap
\begin{probbox}{39}{\stars{1}}
Một đội sửa đường: tuần thứ nhất sửa $\dfrac14$ quãng đường và thêm $3$ km; tuần thứ hai sửa $\dfrac13$ phần còn lại và thêm $1$ km; tuần thứ ba sửa nốt $9$ km cuối. Hỏi quãng đường dài bao nhiêu ki-lô-mét?
\end{probbox}
\probsol{
Sau tuần thứ nhất, $9+1=10$ km là $\dfrac23$ phần còn lại, nên phần còn lại là $15$ km.\par $15+3=18$ km là $\dfrac34$ quãng đường.\par Quãng đường: $18:\dfrac34=24$ km.
\dapso{$24$ km.}}
\gap
\begin{probbox}{40}{\stars{1}}
Cô Hạnh dùng $\dfrac25$ tiền lương cho ăn uống, $\dfrac13$ số còn lại cho việc học của con và gửi tiết kiệm một nửa số còn lại, thì còn giữ $1\,200\,000$ đồng. Hỏi tiền lương của cô là bao nhiêu?
\end{probbox}
\probsol{
Trước khi gửi tiết kiệm: $1\,200\,000\times2=2\,400\,000$ đồng.\par Trước khi chi cho việc học: $2\,400\,000:\dfrac23=3\,600\,000$ đồng.\par Tiền lương: $3\,600\,000:\dfrac35=6\,000\,000$ đồng.
\dapso{$6\,000\,000$ đồng.}}
\gap
\begin{probbox}{41}{\stars{1}}
Lấy tuổi của ông trừ đi $4$, chia cho $7$, cộng thêm $2$, rồi nhân với $5$ thì được $60$. Hỏi ông bao nhiêu tuổi?
\end{probbox}
\probsol{
Trước khi $\times5$: $12$.\par Trước khi $+2$: $10$.\par Trước khi $:7$: $70$.\par Tuổi ông: $70+4=74$.
\dapso{$74$ tuổi.}}
\gap
\begin{probbox}{42}{\stars{1}}
Bèo tấm mỗi ngày tăng diện tích gấp đôi và phủ kín cả ao vào cuối ngày thứ $10$. Hỏi vào ngày thứ mấy bèo phủ được $\dfrac14$ mặt ao?
\end{probbox}
\probsol{
Cuối ngày thứ $9$ bèo phủ $\dfrac12$ mặt ao (gấp đôi lên thì phủ kín).\par Cuối ngày thứ $8$ bèo phủ $\dfrac14$ mặt ao.
\dapso{Ngày thứ $8$.}}
\gap
\begin{probbox}{43}{\stars{1}}
Ba thùng chứa $180$ lít. Rót $10$ l từ thùng I sang thùng II, rồi $20$ l từ thùng II sang thùng III, rồi $5$ l từ thùng III sang thùng I. Lúc này ba thùng bằng nhau. Hỏi lúc đầu mỗi thùng chứa bao nhiêu lít?
\end{probbox}
\probsol{
Lúc cuối mỗi thùng có $180:3=60$ l.\par Trước lần rót thứ ba: I $55$, III $65$.\par Trước lần rót thứ hai: II $80$, III $45$.\par Lúc đầu: I $65$, II $70$, III $45$.
\dapso{I $65$ l, II $70$ l, III $45$ l.}}
\gap
\begin{probbox}{44}{\stars{1}}
Tìm $x$, biết $(x\times1{,}5-2{,}25):0{,}5=7{,}5$.
\end{probbox}
\probsol{
$x\times1{,}5-2{,}25=7{,}5\times0{,}5=3{,}75$.\par $x\times1{,}5=6$.\par $x=4$.
\dapso{$x=4$.}}
\gap
\begin{probbox}{45}{\stars{1}}
Bà tiêu một nửa số tiền và thêm $3\,000$ đồng, rồi một nửa số còn lại và thêm $2\,000$ đồng, rồi một nửa số còn lại và thêm $1\,000$ đồng thì hết tiền. Hỏi bà mang đi bao nhiêu tiền?
\end{probbox}
\probsol{
Trước lần tiêu thứ ba: $(0+1\,000)\times2=2\,000$ đồng.\par Trước lần tiêu thứ hai: $(2\,000+2\,000)\times2=8\,000$ đồng.\par Lúc đầu: $(8\,000+3\,000)\times2=22\,000$ đồng.
\dapso{$22\,000$ đồng.}}
\gap
\begin{probbox}{46}{\stars{1}}
Hai rổ có $60$ quả cam. Chuyển $\dfrac13$ số cam rổ I sang rổ II, rồi chuyển $\dfrac14$ số cam rổ II (lúc đó) sang rổ I. Lúc này số cam hai rổ bằng nhau. Hỏi lúc đầu mỗi rổ có bao nhiêu quả?
\end{probbox}
\probsol{
Lúc cuối mỗi rổ có $30$ quả.\par Trước lần chuyển thứ hai, rổ II có $30:\dfrac34=40$, nên rổ I có $20$.\par Trước lần chuyển thứ nhất, rổ I có $20:\dfrac23=30$, nên rổ II có $30$.
\dapso{Mỗi rổ $30$ quả.}}
\gap
\begin{probbox}{47}{\stars{1}}
Hà làm toán trong $45$ phút, nghỉ $15$ phút, rồi đọc sách trong $1$ giờ $20$ phút, xong lúc $21{:}05$. Hỏi Hà bắt đầu lúc mấy giờ?
\end{probbox}
\probsol{
Bắt đầu đọc sách lúc $21{:}05-1{:}20=19{:}45$.\par Bắt đầu nghỉ lúc $19{:}30$.\par Bắt đầu làm toán lúc $18{:}45$.
\dapso{$18{:}45$.}}
\gap
\begin{probbox}{48}{\stars{1}}
Nghĩ ra một số bất kì, nhân với $9$, cộng thêm $18$, chia cho $9$, rồi trừ đi số đã nghĩ. Chứng tỏ rằng kết quả luôn bằng $2$.
\end{probbox}
\probsol{
Gọi số đã nghĩ là $a$.\par $a\times9+18=(a+2)\times9$.\par Chia cho $9$: $a+2$.\par Trừ đi $a$: $2$, dù $a$ là số nào.
\dapso{Luôn bằng $2$.}}
\gap
\begin{probbox}{49}{\stars{2}}
Một con ốc sên ở đáy một cái giếng sâu $10$ m, ban ngày leo lên được $3$ m, ban đêm tụt xuống $2$ m. Hỏi đến ngày thứ mấy ốc sên lên tới miệng giếng?
\end{probbox}
\probsol{
Tính ngược từ ngày cuối: ốc sên leo $3$ m thì lên tới miệng giếng, nên đầu ngày đó nó ở độ cao $10-3=7$ m.\par Mỗi ngày đêm ốc sên lên được $3-2=1$ m.\par Để lên tới $7$ m cần $7$ ngày đêm.\par Vậy ốc sên ra khỏi giếng vào ngày thứ $8$.
\dapso{Ngày thứ $8$.}}
\gap
\begin{probbox}{50}{\stars{2}}
Bà có một số quả táo. Người con cả lấy một nửa số táo và thêm nửa quả; người con thứ hai lấy một nửa số còn lại và thêm nửa quả; người con út lấy một nửa số còn lại và thêm nửa quả thì vừa hết. Không quả táo nào bị cắt. Hỏi lúc đầu bà có bao nhiêu quả táo?
\end{probbox}
\probsol{
Trước khi người con út lấy: $\left(0+\dfrac12\right)\times2=1$.\par Trước khi người con thứ hai lấy: $\left(1+\dfrac12\right)\times2=3$.\par Lúc đầu: $\left(3+\dfrac12\right)\times2=7$.\par Thử lại: các con lấy $4$, $2$ và $1$ quả táo nguyên.
\dapso{$7$ quả táo.}}
\clearpage
\titlehead
